\Article{20}
It follows readily from elimination theory that the unavoidable errors in observation make the calculation less certain. The more coefficients we must determine by elimination, the greater the uncertainty. This makes the method prescribed in Article $18$, II. particularly valuable, since it suppresses the coefficients of the powers $R^{-(n+2)},R^{-(n+4)}$. With perfect symmetry, these terms would drop out by themselves, but it would be unsafe to rely on that symmetry. A small departure from symmetry would be far less important in the first method than in the second. In the first method, at least, care should be taken with the point $h'$. If this point, from which the distances are measured, lies sufficiently accurately on the magnetic meridian through $h$, there will scarcely be an appreciable difference between $u'-u''$ and $u'''-u''''$. The second method is different, especially if the apparatus requires an eccentric suspension. With this method you will always achieve much less precision whenever there is no room for observations from both sides. The first method is preferable for a further reason. Since in nature $n=2$, it yields twice the value of $L$ that the second method does. If, however, in the second method with eccentric suspension we wish to eliminate the term depending on $R^{-(n+2)}$ as far as possible, \SourceLeaf{113}the point $h'$ must be chosen so that the center of the needle (for $u=u^0$) lies midway between $h$ and $h'$. The calculation that gave this result must, however, be omitted here for brevity.

